% Diagrama de las dos correspondencias de la pizarra de la clase 4, P13.
\begin{tikzpicture}[x=1cm,y=1cm,every node/.style={font=\small}]
\begin{scope}
\node (d1) at (0,3) {$1$};
\node (d2) at (-1,2) {$2$};
\node (d3) at (1,2) {$3$};
\node (d4) at (-1,1) {$4$};
\node (d6) at (1,1) {$6$};
\node (d12) at (0,0) {$12$};
\draw (d1)--(d2)--(d4)--(d12)--(d6)--(d3)--(d1);
\draw (d2)--(d6);
\node at (0,3.65) {Divisores de $12$};
\end{scope}
\draw[<->,draw=rojoDiagramaCIMAT,thick] (1.8,1.5)--(3.2,1.5);
\begin{scope}[xshift=5cm]
\node (h1) at (0,3) {$\ang{[1]}=\mathbb Z/12\mathbb Z$};
\node (h2) at (-1,2) {$\ang{[2]}$};
\node (h3) at (1,2) {$\ang{[3]}$};
\node (h4) at (-1,1) {$\ang{[4]}$};
\node (h6) at (1,1) {$\ang{[6]}$};
\node (h12) at (0,0) {$\ang{[12]}=\{[0]\}$};
\draw (h1)--(h2)--(h4)--(h12)--(h6)--(h3)--(h1);
\draw (h2)--(h6);
\node at (0,3.65) {Subgrupos};
\end{scope}
\end{tikzpicture}
